\section{Suma geométrica y alcance probabilístico}
Un tipo admite a lo sumo $n^D$ colocaciones etiquetadas, fijando sus coordenadas
libres. En la subprobabilidad de llegada se prescriben sus $h$ índices RC por
$C_L x^h e^{O_L(v^2/j)}$ para $v\le L\log n$. El factor exponencial es acotado
pues $j\ge K$. Por las dos partes precedentes y
$n^D x^h j^{-p}=n^{-1}x^{d-1}j^{-b}$, la aportación de una clase de parámetros
está acotada por
\[
 \frac{C_L}{n}\,C^v\,\frac{v^{c+1}}{c!}(4v^2)^{L_s}
            x^{d-1}j^{-b}
 \le\frac{C_L}{n}\, C_1^v v^5 x^{d-1}(v^2/j)^b .
\]
Aquí $v^c/c!\le e^v$, $L_s\le b+2$ y $L_s=O(v)$. Hay sólo un número polinómico
de clases de parámetros enteros por $v$. Primero fijamos la clase de referencia,
$\rho_0\ge2$ permitido y $L>4/\log\rho_0+1$. Después escogemos $\delta_*$ tal
que $C_1\delta_*^{1/3}<1$, reduciéndola también para la excepción siguiente.
Como $v^2/j\le\delta_*$ para $n$ grande y $x\le\delta_*$, obtenemos en grado $v$
\[
 \EE_*\!\left[\sum_{|U|=v}|[U]V|;\mathcal D\right]
 \le \frac{C_L}{n}C_1^v v^{a}\delta_*^{(v-4)/3}.
\]
Su cola desde cualquier grado que tienda a infinito es $o(1/n)$ hasta $L\log n$.
Los grados hasta cuatro aportan $O(1/n)$ por la cota archivada de grado fijo.
\paragraph{Grados superiores y excepción.}
Sea $D'=\mathcal D\cap\{\rho_0^2(\epsilon_A+1/j)\le\eta\}$ y
$N_n=\lceil4\log n/\log\rho_0\rceil$. En $D'$ la cola no colineal desde $N_n+1$
está acotada por $Cj^2\rho_0^{-N_n}=O(n^{-2})$ por las normas de activación
de $F$ y $\Lambda$. En $\mathcal D\setminus D'$ fijamos $N_0=8$ y usamos
$|V-P_8|\le Cn^{28}$ de la parte 06 de \emph{Grados intermedios}.
En esa excepción alguna diagonal tiene más de $R-1$ sitios, con
$R=\lfloor\eta j/(2M_0\rho_0^2)\rfloor$. Si $m=\lceil B\log n\rceil$,
el momento factorial y la unión sobre las líneas dan para la masa de la excepción
\[
 4n C_B e^{O_B(m^2/j)}
       \left(\frac{j^2/n}{R-m+1}\right)^m .
\]
Cada lista de $m$ sitios en una diagonal prescribe $2m$ índices distintos.
Ajustamos la constante de prescripción pues $\log n\le2\log j$.
Escogiendo $\delta_*$ menor que $\eta/(8M_0\rho_0^2)$, la razón es menor que
$1/2$ para $n$ grande, uniformemente. Elegir $B$ fijo grande da $O(n^{-31})$.
Se suma la cota del corte fijo en $\mathcal D$; esto prueba la tesis de la parte 01.
\paragraph{Primera desigualdad restante para el valor aditivo en esta banda.}
Con $t=j/K$, el aporte regular es $O(1/n)$ y el de $P_8$ en la excepción tiene
la cota archivada de grado fijo en la llegada cuando se usa esa entrada fuera de
$\mathcal D$. Con esa misma entrada la condición restante es
\[
 \lim_{T\to\infty}\limsup_{n\to\infty}\
 \sup_{TK\le j\le\min(M,\delta_*n)}
 \frac n{\log(j/K)}
 \left|\EE_*[V-P_8;\mathcal D^c]\right|=0 .
\]
Se entiende esta fórmula donde los valores en la excepción están definidos conforme
al convenio de la llegada. Si sólo se dispone del corte fijo dentro de $\mathcal D$,
se conserva $V$ entero en la esperanza excepcional. Un control suficiente es una
cota conjunta de probabilidad y tamaño del integrando que dé $O(1/n)$ allí.
La escala objetivo es de la media con signo. El criterio terminal completo contiene
$\exp(g_u)-1$ y la contribución de los demás sectores. Esta nota controla el aporte
aditivo regular $V$; sigue abierta la velocidad del error de conteo $R_Q$.
